\section{Endogenous private adoption}\label{sec:short-adoption}
The strategic adoption margin must be checked before evaluating government payoffs. Under complete marginal-cost relief ($d=0$), let
\begin{equation}\label{eq:short-thresholds}
T_A=P-B,\quad T_U=A-C,\quad T_W=A-S,\quad S=\frac{(1-2c)^2}{16}.
\end{equation}
The outsider gains $2T_A-F$ by supporting a regional union's standard because one investment covers two markets; a member's reverse adoption affects only the outsider's market. Under $SW$, foreign firms face either $T_W-F$ or $T_A-F$. Complete best-response calculations in Online Supplement C imply the following sufficient selection-free condition:
\begin{equation}\label{eq:short-F-region}
F_L=\max\{T_W,T_A\}<F<2T_A=F^*.
\end{equation}
In this interval the outsider uniquely adopts the bloc standard in every $SU$, members uniquely refrain from reciprocal adoption, and $SW$ generates no private adoption. For $F>F^*$ there is no adoption under either $SU$ or $SW$. These are sufficient statements on the canonical domain, not a priority claim about asymmetric adoption.

For incomplete cost relief, define the outsider's updated per-member-market gain $T_O(d)=B_R(d)-B$, where $B_R(d)$ is solved below. Since reverse adoptions must also be reoptimized, let $T_W(d)$ and $T_A^R(d)$ denote the exact largest unilateral $SW$ adoption gain and the reverse-member gain conditional on another member's adoption, respectively; their analytical expressions and all other deviation cases appear in Online Supplement D. With
\begin{equation}\label{eq:short-F-residual}
F_L(d)=\max\{0,T_W(d),T_A^R(d)\},\quad F^*(d)=2T_O(d),
\end{equation}
the strict interval $F_L(d)<F<F^*(d)$ likewise selects outsider-only $SU$ adoption and no $SW$ adoption. Its validity relies on the finite full adoption-game deviation audit, not on imposing an arbitrary equilibrium selection. A nonempty interval is exhibited in Proposition~\ref{prop:incomplete-adaptation-open-set}.
